\documentclass[a4paper]{arrowhead}

\usepackage[yyyymmdd]{datetime}
\usepackage{etoolbox}
\usepackage[T1]{fontenc}
\usepackage[utf8]{inputenc}
\usepackage{multirow}
\usepackage{longtable}
\usepackage{amssymb}
\usepackage{float}
\usepackage{tikz}
\usetikzlibrary{positioning,arrows.meta,fit,calc,shapes.misc,shapes.geometric}

\renewcommand{\dateseparator}{-}
\setlength{\emergencystretch}{3em}
\lstset{basicstyle=\ttfamily\footnotesize,breaklines=true,columns=fullflexible,frame=single,rulecolor=\color{gray!50},backgroundcolor=\color{gray!6}}
\tikzset{
  ahsys/.style={draw=ArrowheadBlue,very thick,fill=ArrowheadBlue!6,rounded corners=2pt,minimum width=3.2cm,minimum height=1.1cm,align=center,font=\small\bfseries},
  ahext/.style={draw=ArrowheadGrey,thick,dashed,fill=gray!5,rounded corners=2pt,minimum width=3.2cm,minimum height=1.0cm,align=center,font=\small},
  ahsrv/.style={draw=ArrowheadPurple,thick,fill=ArrowheadPurple!8,rounded corners=6pt,minimum height=0.55cm,inner xsep=4pt,align=center,font=\scriptsize\ttfamily},
  ahcsm/.style={draw=ArrowheadTeal,thick,fill=ArrowheadTeal!8,rounded corners=6pt,minimum height=0.55cm,inner xsep=4pt,align=center,font=\scriptsize\ttfamily},
  ahdb/.style={draw=ArrowheadOrange,thick,fill=ArrowheadOrange!8,cylinder,shape border rotate=90,aspect=0.25,minimum height=0.9cm,align=center,font=\scriptsize},
  ahflow/.style={-{Stealth[length=2.2mm]},thick,ArrowheadGrey},
  ahlabel/.style={font=\scriptsize,fill=white,inner sep=1pt,align=center},
}

%% Special references
\newcommand{\mref}[1]{{\textcolor{ArrowheadPurple}{\hyperref[sec:model:#1]{#1}}}}
\newcommand{\oref}[1]{{\textcolor{ArrowheadBlue}{\hyperref[sec:operations:#1]{#1}}}}
\newcommand{\pdef}[1]{{\textcolor{ArrowheadGrey}{#1\label{sec:model:primitives:#1}\label{sec:model:primitives:#1s}\label{sec:model:primitives:#1es}}}}
\newcommand{\pref}[1]{{\textcolor{ArrowheadGrey}{\hyperref[sec:model:primitives:#1]{#1}}}}

\newrobustcmd\osubsection[3]{
  \addtocounter{subsection}{1}
  \addcontentsline{toc}{subsection}{\protect\numberline{\thesubsection}operation \textcolor{ArrowheadBlue}{#1}}
  \subsection*{
    \thesubsection\quad
    operation
    \textcolor{ArrowheadBlue}{#1}
    (\notblank{#2}{\mref{#2}}{})
    \notblank{#3}{: \mref{#3}}{}
  }
  \label{sec:operations:#1}
}
\newrobustcmd\msubsection[2]{
  \addtocounter{subsection}{1}
  \addcontentsline{toc}{subsection}{\protect\numberline{\thesubsection}#1 \textcolor{ArrowheadPurple}{#2}}
  \subsection*{\thesubsection\quad#1 \textcolor{ArrowheadPurple}{#2}}
  \label{sec:model:#2} \label{sec:model:#2s} \label{sec:model:#2es}
}

\begin{document}

%% Arrowhead Document Properties
\ArrowheadTitle{workflowExecution}
\ArrowheadServiceID{workflowExecution}
\ArrowheadType{Service Description}
\ArrowheadTypeShort{SD}
\ArrowheadVersion{5.4.0}
\ArrowheadDate{2026-11-24}
\ArrowheadAuthor{Documentation maintenance}
\ArrowheadStatus{REVIEW DRAFT}
\ArrowheadContact{Repository maintainers}
\ArrowheadFooter{\href{www.arrowhead.eu}{www.arrowhead.eu}}
\ArrowheadSetup
%%

%% Front Page
\begin{center}
  \vspace*{1cm}
  \huge{\arrowtitle}

  \vspace*{0.2cm}
  \LARGE{\arrowtype}
  \vspace*{1cm}

  \Large{Service ID: \textit{``\arrowid''}}
  \vspace*{\fill}

  

  \vspace*{1cm}
  \vspace*{\fill}

  % Front Page Abstract
  \begin{abstract}
    This document describes the \texttt{workflow\allowbreak{}Execution} service of the Eclipse Arrowhead workflow support systems, which runs work items on behalf of a WorkflowManager. It is provided by the WorkflowExecutor system, service version 1.0.0, documentation edition 5.4.0.
  \end{abstract}

  \vspace*{1cm}
\end{center}

\newpage
%%

%% Table of Contents
\tableofcontents
\newpage
%%

\section{Overview}
\label{sec:overview}
The \texttt{workflow\allowbreak{}Execution} service is provided by every WorkflowExecutor. In push mode the
WorkflowManager selects an executor for an enabled step and hands it the work item with \oref{start}. The executor
runs the execution primitive of the model, $compile \rightarrow O \rightarrow consume$~\cite{varga2026}, with itself
as the consumer: it reserves or renews a lease when the step needs one, has the provider orchestrated for itself,
composes the input, invokes the provider and reports through \texttt{workflow\allowbreak{}Work\allowbreak{}Item}. The service also tells the
WorkflowManager what the executor can do (\oref{service-info}).

The rest of this document is organized as follows.
In Section~\ref{sec:operations}, we describe the abstract message operations provided by the service.
In Section~\ref{sec:model}, we end the document by presenting the data types used by the mentioned operations.

\subsection{Significant Prior Art}
The executor contract of the Arrowhead Choreographer (\texttt{service-info}, \texttt{start},
\texttt{abort}), unchanged since 4.6; Derhamy's HATEOAS workflow executor~\cite{derhamy2018workflow}.

\subsection{How This Service Is Meant to Be Used}
An executor registers \texttt{workflow\allowbreak{}Execution} in the ServiceRegistry with the metadata
\texttt{workflow\allowbreak{}Capabilities}, \texttt{workflow\allowbreak{}Services}, \texttt{workflow\allowbreak{}Workstation} and \texttt{workflowMode}; the
WorkflowManager discovers it there. For each enabled step the WorkflowManager calls \oref{start} on an executor whose
capabilities satisfy $\eta(s)$ and whose services cover $\theta(s)$, preferring the least-loaded one. \oref{start}
returns at once; the result arrives later as a report. The WorkflowManager calls \oref{abort} when a session is
aborted or a step times out.

\subsection{Important Delimitations and Dependencies}
\label{sec:delimitations}
\begin{itemize}
\item Only the configured WorkflowManager (\texttt{workflow.\allowbreak{}manager.\allowbreak{}system-\allowbreak{}name}) or a sysop may start or abort work items (403 otherwise).
\item The executor keeps no durable state; after a restart \oref{get-status} answers \texttt{UNKNOWN} and the step timeout applies when the step's policy sets one. Without it the WorkflowManager keeps waiting for the report.
\item Equipment is invoked through adapters; the reference implementation has the generic HTTP/HTTPS JSON adapter only.
\item An executor never consumes services outside its local cloud unless the step's policy allows it (\texttt{local\allowbreak{}Cloud\allowbreak{}Only = false}).
\end{itemize}

\paragraph{Authorization.} Each operation is authorized separately; \texttt{start} and \texttt{abort} additionally require the caller to be the configured WorkflowManager or a sysop.

\newpage
\section{Service Interface}
\label{sec:operations}
This section describes the interface of the \texttt{workflow\allowbreak{}Execution} service. Each subsection names an abstract operation,
an input type and an output type, in that order. The input type is named inside parentheses, while the output type is
preceded by a colon. Input and output types are only denoted when accepted or returned, respectively, by the operation
in question. All abstract data types named in this section are defined in Section~\ref{sec:model}. The concrete
interface (HTTP/TLS/JSON) is specified in the IDD of the same name.

\begin{figure}[ht!]
  \centering
  \resizebox{0.92\textwidth}{!}{
\begin{tikzpicture}
\node[ahsys] (prov) at (4.6,0) {WorkflowExecutor\\[-1pt]{\normalfont\scriptsize $\ll$provider$\gg$}};
\node[ahsrv] (srv) at (1.6,0) {workflowExecution};
\draw[thick,ArrowheadPurple] (srv.east) -- (prov.west);
\node[ahext] (c0) at (-3.0,0.00) {WorkflowManager\\[-1pt]{\scriptsize $\ll$consumer$\gg$}};
\draw[ahflow] (c0.east) -- (srv.west);
\end{tikzpicture}
  }
  \caption{SysML block view of the \texttt{workflowExecution} service: provider and typical consumers.}
  \label{fig:SD}
\end{figure}

The following interface operations are available.

{\small
\begin{longtable}{|p{4.4cm}|p{5.6cm}|p{5.8cm}|}

\hline
\rowcolor{gray!33} Operation & Input & Output \\ \hline
\endhead
\oref{service-info} & -- & \mref{ExecutorInfo} \\ \hline
\oref{start} & \mref{WorkItem} & -- \\ \hline
\oref{abort} & \mref{AbortRequest} & -- \\ \hline
\oref{get-status} & \mref{WorkItemRef} & \mref{ExecutionState} \\ \hline
\end{longtable}
}

\osubsection{service-info}{}{ExecutorInfo}
Returns the executor's name, capabilities $cap(w)$, runnable services, workstation and mode.
\osubsection{start}{WorkItem}{}
Accepts a work item and runs it asynchronously. The outcome is reported through \texttt{workflow\allowbreak{}Work\allowbreak{}Item.\allowbreak{}report}.
\osubsection{abort}{AbortRequest}{}
Aborts a running work item; the executor reports ABORTED.
\osubsection{get-status}{WorkItemRef}{ExecutionState}
Returns the executor's local state of a work item.

\newpage
\section{Information Model}
\label{sec:model}
Here, all data objects that can be part of the \texttt{workflow\allowbreak{}Execution} service are listed in alphabetic order. Each
subsection, which describes one type of object, begins with the \textit{struct} keyword (a collection of named fields,
each with its own data type), the \textit{enum} keyword (a closed set of values) or the \textit{alias} keyword. As a
complement to the explicitly defined types in this section, there is also a list of implicit primitive types in
Section~\ref{sec:model:primitives}.

\msubsection{struct}{AbortRequest}
Aborts a work item on an executor.

{\small
\begin{longtable}{|p{3.9cm}|p{4.4cm}|p{7.5cm}|}

\hline
\rowcolor{gray!33} Field & Type & Description \\ \hline
\endhead
\texttt{workItemId} & \pref{WorkItemId} & The work item. \\ \hline
\texttt{sessionId} & \pref{SessionId} & Its session. \\ \hline
\texttt{reason} & \pref{String} & Why. Optional. \\ \hline
\end{longtable}
}

\msubsection{enum}{DispatchMode}
How a step reaches an executor.

{\small
\begin{longtable}{|p{3.9cm}|p{11.9cm}|}

\hline
\rowcolor{gray!33} Value & Meaning \\ \hline
\endhead
\texttt{push} & The WorkflowManager selects an executor and calls \texttt{workflow\allowbreak{}Execution.\allowbreak{}start} (default). \\ \hline
\texttt{pull} & The WorkflowManager publishes the work item; an executor claims it with \texttt{workflow\allowbreak{}Work\allowbreak{}Item.\allowbreak{}claim}. \\ \hline
\texttt{event} & Reserved for EventHandler notification; currently handled as \texttt{pull}. \\ \hline
\end{longtable}
}

\msubsection{struct}{ExecutionState}
State of a work item on an executor.

{\small
\begin{longtable}{|p{3.9cm}|p{4.4cm}|p{7.5cm}|}

\hline
\rowcolor{gray!33} Field & Type & Description \\ \hline
\endhead
\texttt{workItemId} & \pref{WorkItemId} & The work item. \\ \hline
\texttt{state} & \mref{ExecutorState} & Local state. \\ \hline
\texttt{message} & \pref{String} & Last message. Optional. \\ \hline
\end{longtable}
}

\msubsection{struct}{ExecutorInfo}
What an executor declares about itself.

{\small
\begin{longtable}{|p{3.9cm}|p{4.4cm}|p{7.5cm}|}

\hline
\rowcolor{gray!33} Field & Type & Description \\ \hline
\endhead
\texttt{name} & \pref{SystemName} & Executor system name. \\ \hline
\texttt{capabilities} & \pref{List}$<$\pref{String}$>$ & Capabilities $cap(w)$. \\ \hline
\texttt{services} & \pref{List}$<$\pref{ServiceDefinition}$>$ & Services it can run; \texttt{*} means any. \\ \hline
\texttt{workstation} & \pref{String} & Workstation. Optional. \\ \hline
\texttt{mode} & \mref{DispatchMode} & \texttt{push} or \texttt{pull}. \\ \hline
\end{longtable}
}

\msubsection{struct}{ExecutorRequirement}
$\eta(s)$: constraints on the executor that runs a step.

{\small
\begin{longtable}{|p{3.9cm}|p{4.4cm}|p{7.5cm}|}

\hline
\rowcolor{gray!33} Field & Type & Description \\ \hline
\endhead
\texttt{capability} & \pref{String} & A capability the executor must declare in \texttt{workflow\allowbreak{}Capabilities}. Optional. \\ \hline
\texttt{workstation} & \pref{String} & The executor's workstation. Optional. \\ \hline
\texttt{system} & \pref{SystemName} & A named executor system. Optional. \\ \hline
\end{longtable}
}

\msubsection{enum}{ExecutorState}
Local state of a work item on an executor. The authoritative state is held by the WorkflowManager.

{\small
\begin{longtable}{|p{3.9cm}|p{11.9cm}|}

\hline
\rowcolor{gray!33} Value & Meaning \\ \hline
\endhead
\texttt{ACCEPTED} & Received, not yet started. \\ \hline
\texttt{RUNNING} & Being run. \\ \hline
\texttt{DONE} & Run finished and reported (whatever the report was). \\ \hline
\texttt{UNKNOWN} & The executor does not know the work item, e.g. after a restart. \\ \hline
\end{longtable}
}

\msubsection{enum}{LeaseMode}
Role of a step for a shared resource.

{\small
\begin{longtable}{|p{3.9cm}|p{11.9cm}|}

\hline
\rowcolor{gray!33} Value & Meaning \\ \hline
\endhead
\texttt{reserve} & Reserves the resource; its result carries the lease token. \\ \hline
\texttt{use} & Uses the resource; the executor renews the lease immediately before the invocation. \\ \hline
\texttt{release} & Releases the resource. \\ \hline
\end{longtable}
}

\msubsection{struct}{LeasePolicy}
The role of a step for a shared resource.

{\small
\begin{longtable}{|p{3.9cm}|p{4.4cm}|p{7.5cm}|}

\hline
\rowcolor{gray!33} Field & Type & Description \\ \hline
\endhead
\texttt{resource} & \pref{Name} & A declared resource. \\ \hline
\texttt{mode} & \mref{LeaseMode} & \texttt{reserve}, \texttt{use} or \texttt{release}. \\ \hline
\end{longtable}
}

\msubsection{enum}{OnError}
Reaction to an \texttt{ERROR} report.

{\small
\begin{longtable}{|p{3.9cm}|p{11.9cm}|}

\hline
\rowcolor{gray!33} Value & Meaning \\ \hline
\endhead
\texttt{rebind} & Exclude the executor for this session and bind the step again. \\ \hline
\texttt{retry} & Keep the executor: after \texttt{backoff\allowbreak{}Seconds} dispatch the step to the same executor again, or to the least-loaded candidate when that executor is no longer available (extension x-retry; for transient provider-side failures, the provider must accept a repeated call). \\ \hline
\texttt{abort} & Fail the session. \\ \hline
\end{longtable}
}

\msubsection{struct}{Policy}
Execution policy of a step. Every field other than \texttt{lease} and \texttt{local\allowbreak{}Cloud\allowbreak{}Only} is an extension outside the proven model.

{\small
\begin{longtable}{|p{3.9cm}|p{4.4cm}|p{7.5cm}|}

\hline
\rowcolor{gray!33} Field & Type & Description \\ \hline
\endhead
\texttt{quantity} & \pref{Integer} & Number of executions; only 1 is accepted (check M01). Optional. \\ \hline
\texttt{timeout\allowbreak{}Seconds} & \pref{Duration} & Extension x-timeout: the step fails when it runs longer. Optional. \\ \hline
\texttt{nominal\allowbreak{}Duration\allowbreak{}Seconds} & \pref{Duration} & Expected duration, used to check lease durations (check L04). Optional. \\ \hline
\texttt{retry} & \mref{Retry} & Extensions x-rebind and x-retry: behaviour on \texttt{ERROR}. Optional. \\ \hline
\texttt{lease} & \mref{LeasePolicy} & The step's role for a shared resource. Optional. \\ \hline
\texttt{local\allowbreak{}Cloud\allowbreak{}Only} & \pref{Boolean} & Whether orchestration may return providers of other clouds (default true). Optional. \\ \hline
\end{longtable}
}

\msubsection{struct}{Retry}
Extensions x-rebind and x-retry.

{\small
\begin{longtable}{|p{3.9cm}|p{4.4cm}|p{7.5cm}|}

\hline
\rowcolor{gray!33} Field & Type & Description \\ \hline
\endhead
\texttt{maxAttempts} & \pref{Integer} & Maximum attempts of the step in one session, counting every work item of the step; at least 1, default 1 (no further attempt). Optional. \\ \hline
\texttt{onError} & \mref{OnError} & What an \texttt{ERROR} report causes while attempts remain. State it explicitly: the recipe schema names \texttt{rebind} as default, but the Java WorkflowManager fails the session when it is absent. Optional. \\ \hline
\texttt{backoff\allowbreak{}Seconds} & \pref{Duration} & Pause before a same-executor retry (\texttt{onError} = \texttt{retry}); 0--600, default 1. Optional. \\ \hline
\end{longtable}
}

\msubsection{struct}{ServiceRequirement}
Requirements on the provider of the step's service, passed to orchestration.

{\small
\begin{longtable}{|p{3.9cm}|p{4.4cm}|p{7.5cm}|}

\hline
\rowcolor{gray!33} Field & Type & Description \\ \hline
\endhead
\texttt{service\allowbreak{}Definition} & \pref{ServiceDefinition} & The service to consume. \\ \hline
\texttt{interfaces} & \pref{List}$<$\pref{String}$>$ & Accepted interface template names, e.g. \texttt{generic\_\allowbreak{}https}. Optional. \\ \hline
\texttt{minVersion} & \pref{String} & Lowest accepted service version. Optional. \\ \hline
\texttt{maxVersion} & \pref{String} & Highest accepted service version. Optional. \\ \hline
\texttt{metadata} & \pref{Map}$<$\pref{String}, \pref{JsonValue}$>$ & Metadata requirements on the service instance. Optional. \\ \hline
\texttt{qos} & \pref{Map}$<$\pref{String}, \pref{JsonValue}$>$ & QoS requirements forwarded to orchestration. Optional. \\ \hline
\texttt{requires\allowbreak{}Results} & \pref{List}$<$\pref{Name}$>$ & Steps whose results the provider needs; each must be in $dep(s)$ (check D04). Optional. \\ \hline
\end{longtable}
}

\msubsection{struct}{WorkItem}
One attempt to run one step: what an executor needs to run it.

{\small
\begin{longtable}{|p{3.9cm}|p{4.4cm}|p{7.5cm}|}

\hline
\rowcolor{gray!33} Field & Type & Description \\ \hline
\endhead
\texttt{workItemId} & \pref{WorkItemId} & Identifier of the attempt. \\ \hline
\texttt{sessionId} & \pref{SessionId} & Session. \\ \hline
\texttt{recipe} & \pref{String} & \texttt{name@version} of the recipe. \\ \hline
\texttt{action} & \pref{Name} & Action of the step. \\ \hline
\texttt{step} & \pref{Name} & Step name. \\ \hline
\texttt{attempt} & \pref{Integer} & Attempt number from 1. \\ \hline
\texttt{service} & \mref{ServiceRequirement} & $\theta(s)$. \\ \hline
\texttt{executor} & \mref{ExecutorRequirement} & $\eta(s)$. Optional. \\ \hline
\texttt{params} & \pref{Map}$<$\pref{String}, \pref{JsonValue}$>$ & Static parameters of the step. Optional. \\ \hline
\texttt{input} & \pref{Template} & Input template $\kappa_s$. Optional. \\ \hline
\texttt{policy} & \mref{Policy} & Step policy. Optional. \\ \hline
\texttt{lease\allowbreak{}Duration\allowbreak{}Seconds} & \pref{Duration} & Lease duration of the step's resource. Optional. \\ \hline
\texttt{parameters} & \pref{JsonObject} & Session parameters $\rho$ (included for convenience; also available from \texttt{workflow\allowbreak{}Context}). Optional. \\ \hline
\texttt{dependency\allowbreak{}Results} & \pref{Map}$<$\pref{Name}, \pref{JsonValue}$>$ & $\xi|dep(s)$ (also available from \texttt{workflow\allowbreak{}Context}). Optional. \\ \hline
\end{longtable}
}

\msubsection{struct}{WorkItemRef}
Identifies a work item.

{\small
\begin{longtable}{|p{3.9cm}|p{4.4cm}|p{7.5cm}|}

\hline
\rowcolor{gray!33} Field & Type & Description \\ \hline
\endhead
\texttt{workItemId} & \pref{WorkItemId} & The work item. \\ \hline
\end{longtable}
}


\subsection{Primitives}
\label{sec:model:primitives}
Types and structures mentioned throughout this document that are assumed to be available to implementations of this
service. The concrete interpretation of each of these types must be provided by any IDD document claiming to implement
this service.

{\small
\begin{longtable}{|p{3.2cm}|p{12.6cm}|}

\hline
\rowcolor{gray!33} Type & Description \\ \hline
\endhead
\pdef{Boolean} & One out of \texttt{true} or \texttt{false}. \\ \hline
\pdef{Duration} & A length of time in seconds. \\ \hline
\pdef{Integer} & A whole number. \\ \hline
\pdef{JsonObject} & A set of named values. \\ \hline
\pdef{JsonValue} & Any value: object, array, string, number, Boolean or null. \\ \hline
\pdef{List} & An ordered collection of items of the type given in angle brackets. \\ \hline
\pdef{Map} & A collection of key/value pairs of the types given in angle brackets. \\ \hline
\pdef{Name} & An identifier of a recipe, action, step, parameter or resource: a letter followed by at most 62 letters, digits, \texttt{-} or \texttt{\_}. \\ \hline
\pdef{ServiceDefinition} & The name of an Arrowhead service definition. Names are compared case-insensitively and without \texttt{-}, \texttt{\_} and spaces, so \texttt{temperature-\allowbreak{}read} and \texttt{temperature\allowbreak{}Read} denote the same service. \\ \hline
\pdef{SessionId} & The identifier the WorkflowManager gives a session. \\ \hline
\pdef{String} & A chain of characters. \\ \hline
\pdef{SystemName} & The name of an Arrowhead system as it appears in its certificate and in the ServiceRegistry. \\ \hline
\pdef{Template} & A value in which strings of the form \texttt{\$\{param.$p$\}} or \texttt{\$\{result.$s$.path\}} are replaced by a session parameter or by (a part of) a stored step result. \\ \hline
\pdef{WorkItemId} & The identifier of one attempt to run one step of one session. \\ \hline
\end{longtable}
}


\newpage
\bibliographystyle{IEEEtran}
\bibliography{bibliography}

\newpage

\section{Revision History}
\subsection{Amendments}

\noindent\begin{tabularx}{\textwidth}{| p{1cm} | p{2.2cm} | p{1.6cm} | X | p{3.6cm} |} \hline
\rowcolor{gray!33} No. & Date & Version & Subject of Amendments & Author \\ \hline
1 & 2026-11-23 & \arrowversion & Initial edition: service description of workflowExecution, written from the Java reference implementation (ah5-Workflow-java). & Documentation maintenance \\ \hline
2 & 2026-11-24 & \arrowversion & Retry policy as implemented; behaviour after a restart. & Documentation maintenance \\ \hline
\end{tabularx}

\subsection{Quality Assurance}

\noindent\begin{tabularx}{\textwidth}{| p{1cm} | p{2.2cm} | p{1.6cm} | X |} \hline
\rowcolor{gray!33} No. & Date & Version & Approved by \\ \hline
1 & & \arrowversion & Pending review \\ \hline
\end{tabularx}

\end{document}
